\chapter{The Operator Portal}
\label{ch:operator}

\portalcover{Operator Portal}{For staff working at a station on the factory floor}{secStitch}

\section{Who this chapter is for}

You work at one of the stations on the floor: the fabric store, the cutting
room, a sewing line, the sorting table, the wash plant, finishing, the retrofit
bench or dispatch. Your job in the system is simple:

\begin{plainbox}
\textbf{Scan what arrives. Scan what leaves. Tell the system what you did.}
\end{plainbox}

You do not need to understand the rest of the factory. Each screen tells you what
is waiting for you and what to press.

\section{Your daily routine}

\begin{figure}[H]
\centering
\begin{tikzpicture}[
  node distance=7mm,
  st/.style={rectangle,rounded corners=3pt,fill=brandsoft,draw=brandblue,
             minimum height=13mm,text width=30mm,align=center,font=\small},
  ar/.style={-{Stealth[length=3mm]},thick,draw=brandblue}]
\node[st] (a) {\textbf{1.} Sign in};
\node[st,right=of a] (b) {\textbf{2.} Check what is waiting};
\node[st,right=of b] (c) {\textbf{3.} Receive it and count};
\node[st,right=of c] (d) {\textbf{4.} Do your work};
\node[st,below=8mm of d] (e) {\textbf{5.} Send it onward};
\node[st,left=of e] (f) {\textbf{6.} Print the note};
\node[st,left=of f] (g) {\textbf{7.} Sign out};
\draw[ar] (a)--(b); \draw[ar] (b)--(c); \draw[ar] (c)--(d); \draw[ar] (d)--(e);
\draw[ar] (e)--(f); \draw[ar] (f)--(g);
\end{tikzpicture}
\caption{The same seven steps apply at every station.}
\end{figure}

\section{Moving goods in and out}

Receiving a batch into your section and sending one onward are the same at
every station, and are described in Section~\ref{sec:receive} and
Section~\ref{sec:dispatch}. The rest of this chapter covers the work that is
particular to each station.

\section{Fabric store: receiving denim rolls}

Open \menu{Fabric Store}. The first tab shows how much fabric you are holding.

\shot{30-store-stock}{Fabric stock, grouped by type, colour and condition.}{fig:store-stock}

To book in a delivery, press \btn{+ Receive rolls}.

\shotsmall{32-store-receive}{Booking in a delivery. Add one block per roll.}{fig:store-receive}

\begin{steps}
  \item Fill in the \field{Supplier} and \field{Invoice reference} at the top.
  \item For each roll, choose the fabric type and colour, then enter the shade
        or batch, the length in metres, and where you are putting it.
  \item If the roll carries its own tag, scan it into \field{Roll tag EPC}.
        This lets the cutting room pull rolls by scanning instead of typing.
  \item \btn{+ Duplicate last} copies the previous roll's details --- useful for
        a pallet of the same fabric.
  \item Press \btn{Receive rolls}.
\end{steps}

The system creates the goods receipt number for you. Every roll is listed on the
\menu{Roll register} tab, where you can search by roll number, tag, batch or
location.

\shot{31-store-rolls}{The roll register. Every roll ever received, with how much
is left on it.}{fig:store-rolls}

\subsection{Looking up a past delivery}

The \menu{Goods receipts} tab lists every delivery the store has ever booked in,
newest first, with the supplier, their invoice number, how many rolls came and
how many metres that was --- and who received it.

\shot{33-store-grn-list}{The goods receipt register. Click any row to see the
individual rolls that arrived on that delivery.}{fig:store-grn}

\begin{tipbox}
This is the screen to open when a supplier queries a delivery. The \field{GRN no}
is the factory's own receipt number; the \field{Invoice} column is the
supplier's. Quoting the GRN number settles most questions in one telephone call.
\end{tipbox}

\section{Cutting: issuing fabric and making bundles}

Open \menu{Cutting}. Each line is a cut order --- an instruction to cut a
particular design and colour.

\shot{40-cutting-list}{Cut orders, showing how many pieces have been cut and how
many have been tagged so far.}{fig:cut-list}

\subsection{Starting a new cut order}

A cut order says what is going to be cut, before any fabric is issued. Press
\btn{+ New cut order} at the top right.

\shotsmall{42-cutting-new}{Creating a cut order. Only the design, colour and
planned quantity are required.}{fig:cut-new}

\begin{steps}
  \item Choose the \field{Design / style} --- for example
        \typed{JKT-201 $\cdot$ Trucker Denim Jacket}.
  \item Choose the \field{Colour}.
  \item If this cutting is for a particular customer order, pick it under
        \field{Customer order}. Leave it on \emph{Not linked to an order} for
        stock production.
  \item Type the \field{Planned quantity} --- how many pieces you intend to cut.
  \item Add \field{Remarks} if the marker or lay needs a note.
  \item Press \btn{Create}.
\end{steps}

\begin{tipbox}
Linking the cut order to a customer order is worth doing. It is what lets the
\menu{Orders} screen show that customer how many of their pieces are finished,
without anybody counting anything by hand.
\end{tipbox}

Click a cut order to open it.

\shotsmall{41-cutting-detail}{Inside a cut order: the fabric consumed and the
bundles produced.}{fig:cut-detail}

\subsection{Issuing rolls to a cut}

\begin{steps}
  \item Press \btn{Issue rolls}.
  \item Only rolls of the matching colour are offered, so you cannot issue the
        wrong shade by mistake.
  \item Either tick the rolls, or scan their tags and press
        \btn{Select scanned rolls}.
  \item Adjust the metres if you are not using the whole roll.
  \item Press \btn{Issue selected}.
\end{steps}

\begin{warnbox}
The system will not let you issue more metres than a roll actually has left. If
you see that message, check the roll number --- you may have picked up the wrong one.
\end{warnbox}

\subsection{Creating bundles}

\begin{steps}
  \item Press \btn{Create bundles}.
  \item Add a line for each size: how many bundles, and how many pieces in each.
  \item The running total is shown at the bottom. Check it against your lay.
  \item Press \btn{Create bundles}. Each bundle gets its own number.
\end{steps}

When the bundles are tied and ready, press \btn{Issue all to stitching}.

\section{Stitching: counting in and attaching tags}

This is the most important station in the whole system, because this is where a
garment first gets its identity.

\shot{50-stitching}{The stitching screen. Bundles waiting to be counted are at the
top; bundles being tagged are underneath.}{fig:stitching}

\subsection{Step one: count the bundle in}

The bundle arrives with a paper count from cutting. Your job is to check it.

\begin{steps}
  \item Press \btn{Count in} on the bundle.
  \item Physically count the pieces.
  \item Type the number you actually counted.
  \item Press \btn{Confirm count}.
\end{steps}

\shotsmall{51-stitching-count}{Counting a bundle in. Type what you really
counted, not what the paper says.}{fig:stitch-count}

\begin{tipbox}
\textbf{Always type your real count.} If it differs from cutting's figure, the
system records the difference and moves on --- nobody gets into trouble, and the
factory learns where pieces are being lost. Hiding a difference is what causes
problems later.
\end{tipbox}

\subsection{Step two: attach the tags}

\begin{steps}
  \item Press \btn{Attach tags} on the bundle.
  \item Fit a tag to each finished garment.
  \item Scan all the tags into the box.
  \item Press \btn{Register garments}.
\end{steps}

\shotsmall{52-stitching-tags}{Attaching tags. The three boxes at the top show the
bundle quantity, how many are already tagged, and how many are left.}{fig:stitch-tags}

\begin{stopbox}
\textbf{One tag, one garment, forever.} If a tag has already been used on another
garment, the system refuses it and shows you which ones. Never try to force a
duplicate tag through --- it would make two garments look like the same one and
the count would be wrong from then on. Put that tag aside and use a fresh one.
\end{stopbox}

You also cannot register more tags than the bundle has pieces. If you get that
message, re-count the bundle.

\section{Sorting: splitting a pile into groups}

The sorting table is where a mixed pile is separated into batches. The system
does the sorting for you --- you just read the pile.

\shot{60-sorting-list}{Sorting sessions. Open sessions are listed at the top.}{fig:sort-list}

\subsection{Starting a session}

\begin{steps}
  \item Press \btn{+ New sorting session}.
  \item Choose your \field{Station}.
  \item Choose how to group the pile. Two ready-made choices cover nearly
        everything:
        \begin{bullets}
          \item \btn{Before wash: design / colour / size}
          \item \btn{After wash: order / size}
        \end{bullets}
  \item Press \btn{Open session}.
\end{steps}

\shotsmall{61-sorting-new}{Starting a sorting session. The preset buttons set the
grouping for you.}{fig:sort-new}

\subsection{Reading and sending}

Bulk-read the pile into the scan box and press \btn{Add to session}. The system
splits the read into groups and counts each one.

\shotsmall{62-sorting-session}{A sorting session. Each group has its own row and
its own send button.}{fig:sort-session}

Each group has a \btn{Send to \ldots} button. Pressing it creates a separate
transfer note for that group alone, so the wash plant receives one clean batch at
a time.

\begin{warnbox}
Anything the system cannot place appears in an \textbf{Exceptions} list ---
usually a garment from another section, or a tag it does not recognise. Take
those pieces off the table and show them to your supervisor. They are not in
any group and will not be sent.
\end{warnbox}

\section{Retrofit bench: correcting a fault}

When quality control fails a garment, it comes to you with everything the
inspector recorded.

\shot{90-retrofit-queue}{The retrofit queue. The \field{Defects} column shows how
many faults are still outstanding on each garment.}{fig:retro-queue}

\begin{steps}
  \item Scan the garment's tag into the box at the top, or click its row in the
        queue.
  \item The screen fills with the design picture and numbered markers showing
        \emph{exactly} where each fault is.
  \item Make the corrections.
  \item Tick off the faults you have fixed (all are ticked already).
  \item Write what you did in \field{Correction carried out}.
  \item Press \btn{Mark corrected}.
\end{steps}

\shot{91-retrofit-file}{The defect file. The numbered marker on the design shows
precisely where the inspector found the problem.}{fig:retro-file}

\begin{tipbox}
Write the correction in normal words --- ``re-stitched left knee seam and
pressed'' is perfect. The inspector reads this when the garment comes back, so
they know what to look at.
\end{tipbox}

When a batch is corrected, send it back to quality control using
Section~\ref{sec:dispatch}.

\section{Dispatch: changing the tag}

At dispatch the factory's own tracking tag comes off and the customer's tag goes
on.

\shot{95-dispatch-ready}{What is waiting to be re-tagged, grouped by style, colour
and size.}{fig:disp-ready}

\begin{steps}
  \item Open the shipment from the \menu{Shipments} tab, or press
        \btn{Start a shipment}.
  \item Press \btn{Re-tag garments}.
  \item Scan the \textbf{tracking tag} on the garment. Its details appear,
        including the tag specification that this customer requires.
  \item Apply the customer's tag to the garment.
  \item Scan the \textbf{customer tag}. The pair is added to the list.
  \item Repeat for the tray, then press \btn{Apply tag swap}.
\end{steps}

\shotsmall{97-dispatch-swap}{Re-tagging. Scan the old tag, then the new one, and
press Enter after each scan.}{fig:disp-swap}

\begin{plainwords}
After the swap, the garment answers to its \emph{new} customer tag. The old
factory tag is released, so it can be put on a brand new garment tomorrow. If you
scan that old tag now, the system correctly says ``not registered'' --- because
that tag no longer belongs to this pair of jeans.
\end{plainwords}

\subsection{Shipments: packing and sending}

Re-tagged garments go into a \textbf{shipment} --- one carton load for one
customer. The \menu{Shipments} tab lists them.

\shot{96-dispatch-shipments}{The shipment list. \texttt{OPEN} means you can still
add garments to it; \texttt{SHIPPED} means it has left the building.}{fig:disp-shipments}

\begin{steps}
  \item Press \btn{+ New shipment} and choose the customer, the order it is
        against, and the carrier.
  \item The shipment opens with a status of \texttt{OPEN} and zero units.
  \item Click that shipment, then re-tag garments into it as described above.
        The \field{Units} count rises as you go.
  \item When the carton is full, mark the shipment as shipped. The packing list
        prints automatically.
\end{steps}

\begin{warnbox}
Once a shipment is marked \texttt{SHIPPED} you cannot add anything more to it.
Check the \field{Units} count against the customer's order before you close it.
\end{warnbox}
